%% Corrigendum to Livingstone et al. (2019), Bernoulli 25(4A), 3109--3138.
%% IMPORTANT: affiliations and email addresses below are those printed in the
%% original 2019 article. Update them before submission if appropriate.

%\documentclass[bj,authoryear,noshowframe]{imsart}
%\documentclass[bj,numbers,noshowframe]{imsart}

%\usepackage{amssymb}
%\setcitestyle{open={[},close={]}}

\startlocaldefs
\theoremstyle{plain}
\newtheorem*{lemma515}{Lemma 5.15}
\newtheorem*{lemma516}{Lemma 5.16}
\newtheorem*{lemma517}{Lemma 5.17}
\newtheorem*{theorem514}{Theorem 5.14}
\newtheorem*{corollary23}{Corollary 2.3(ii), $\beta>2$ case}
\endlocaldefs

%\begin{document}

\begin{frontmatter}

\title{Corrigendum to ``On the geometric ergodicity of Hamiltonian Monte Carlo''}
\runtitle{Corrigendum: geometric ergodicity of HMC}

\begin{aug}
\author[A]{\inits{S.}\fnms{Samuel}~\snm{Livingstone}\ead[label=e1]{samuel.livingstone@ucl.ac.uk}}
\author[B]{\inits{M.}\fnms{Michael}~\snm{Betancourt}\ead[label=e2]{betan@symplectomorphic.com}}
\author[C]{\inits{S.}\fnms{Simon}~\snm{Byrne}\ead[label=e3]{simonbyrne@gmail.com}}
\author[D]{\inits{M.}\fnms{Mark}~\snm{Girolami}\ead[label=e4]{mag92@cam.ac.uk}}

\address[A]{Department of Statistical Science, University College London,
Gower Street, London WC1E 6BT, United Kingdom\printead[presep={,\ }]{e1}}
\address[B]{Symplectomorphic LLC, New York, NY, USA\printead[presep={,\ }]{e2}}
\address[C]{NVIDIA Corporation, 2788 San Tomas Expressway,
Santa Clara, CA 95051, USA\printead[presep={,\ }]{e3}}
\address[D]{Department of Engineering, University of Cambridge, Trumpington Street, Cambridge CB2 1PZ, United Kingdom\printead[presep={,\ }]{e4}}


\runauthor{S. Livingstone et al.}
\end{aug}

\begin{abstract}
We correct a consequential typographical error in Theorem 5.14 of \cite{Livingstone2019}. Condition (18) should require
$\|y\|\geq 3\|x\|$, rather than $\|y\|\geq 2\|x\|$, while retaining the
requirement $\|\nabla U(y)\|\geq 3\|\nabla U(x)\|$. The original condition is
not satisfied by the exponential family $E(\beta,\alpha)$ for every
$\beta>2$, meaning Corollary 2.3 was incomplete as stated. We give replacement proofs of Lemmas 5.15--5.17 and
Theorem 5.14 for completeness.  For
$E(\beta,\alpha)$ the corrected condition holds for every $\beta>2$, meaning Corollary 2.3 remains valid as stated.
\end{abstract}

\begin{keyword}
\kwd{geometric ergodicity}
\kwd{Hamiltonian Monte Carlo}
\kwd{leapfrog integrator}
\end{keyword}

\end{frontmatter}

\section{Details of the typographical error and correction}

In \cite{Livingstone2019}, p.~3128, condition (18), used in Theorem 5.14 and
then in part (i) of Theorem 2.2, states that there is a fixed $C<\infty$ such
that whenever $\|y\|\geq2\|x\|\geq C$ then
$\|\nabla U(y)\|\geq3\|\nabla U(x)\|$. For the one-dimensional exponential
family $E(\beta,\alpha)$, for which $U(x)=\alpha|x|^\beta+\mathrm{const.}$ in
the tails,
\[
 \frac{\|\nabla U(y)\|}{\|\nabla U(x)\|}
 \geq 2^{\beta-1}
\]
whenever $|y|\geq2|x|$ and both points are in the tails.  The
printed condition is therefore guaranteed only when
$\beta\geq1+\log(3)/\log(2)$, rather than for every $\beta>2$. The proof
of Corollary 2.3(ii) therefore has a gap for
\[
 2<\beta<1+\frac{\log 3}{\log 2}.
\]
Condition (18) should instead read: there is a fixed $C<\infty$ such that
whenever $\|y\|\geq3\|x\|\geq C$ then
\begin{equation}
 \|\nabla U(y)\|\geq3\|\nabla U(x)\|.
 \tag{18}
 \label{eq:18}
\end{equation}
Condition (19) can be left unchanged. With these updates Corollary 2.3(ii) is true as stated, lack of geometric ergodicity is proven for any $\beta > 2$.  The conclusion of Theorem 5.14, and therefore
the conclusion used in Theorem 2.2, is unchanged. We thank the authors of \citep{gruffaz2026theoretical} for bringing the error to our attention.

There is another (albeit inconsequential) typographical error in the printed proof.  Lemma 5.16 is stated under
condition (17) alone, but its proof invokes condition (18). The following
minimally-changed replacement proofs show directly that condition (17) alone is sufficient for the lemma.

\section{Replacement proofs}

Fix a leapfrog step size $\varepsilon>0$. The first leapfrog position update
is
\begin{equation}
 x_{\varepsilon}
 =x_0-\frac{\varepsilon^2}{2}\nabla U(x_0)+\varepsilon p_0.
 \label{eq:firststep}
\end{equation}
For each $i\geq1$, eliminating the intermediate half-step momentum from two
successive leapfrog position updates gives the exact recurrence
\begin{equation}
 x_{(i+1)\varepsilon}
 =2x_{i\varepsilon}-x_{(i-1)\varepsilon}
 -\varepsilon^2\nabla U(x_{i\varepsilon}).
 \label{eq:recurrence}
\end{equation}
We retain condition (17) of \cite{Livingstone2019}, namely
\begin{equation}
 \lim_{\|x\|\to\infty}
 \frac{\|\nabla U(x)\|}{\|x\|}=\infty.
 \tag{17}
 \label{eq:17}
\end{equation}

\begin{lemma515}
If \eqref{eq:17} holds, then there exists an $\eta<\infty$ such that for all
$\|x_0\|>\eta$ and any $\|p_0\|\leq\|x_0\|^\delta$ for any fixed $\delta<1$, it
holds that
\[
 \|x_{\varepsilon}\|>3\|x_0\|.
\]
\end{lemma515}

\begin{proof}
Taking norms in \eqref{eq:firststep} gives
\[
 \|x_{\varepsilon}\|
 \geq
 \frac{\varepsilon^2}{2}\|\nabla U(x_0)\|
 -\|x_0\|-\varepsilon\|p_0\|.
\]
Dividing by $\|x_0\|$ gives
\begin{equation}
 \frac{\|x_{\varepsilon}\|}{\|x_0\|}
 \geq
 \frac{\varepsilon^2}{2}
 \frac{\|\nabla U(x_0)\|}{\|x_0\|}
 -1-\varepsilon\frac{\|p_0\|}{\|x_0\|}.
 \label{eq:firststepbound}
\end{equation}
The first term on the right-hand side tends to infinity by \eqref{eq:17},
while the last term is bounded above by
$\varepsilon\|x_0\|^{\delta-1}$ and tends to zero. Hence the right-hand side
of \eqref{eq:firststepbound} is larger than $3$ for all sufficiently large
$\|x_0\|$.
\end{proof}

\begin{lemma516}
If \eqref{eq:17} holds, then there exists an $\eta<\infty$ such that for all
$\|x_0\|>\eta$ and any $\|p_0\|\leq\|x_0\|^\delta$ for any fixed $\delta<1$, it
holds that
\begin{equation}
 \|x_{i\varepsilon}\|\geq3\|x_{(i-1)\varepsilon}\|,
 \qquad i=1,\ldots,L,
 \label{eq:stepwiseexpansion}
\end{equation}
meaning
\begin{equation}
 \|x_{L\varepsilon}\|\geq3^L\|x_0\|.
 \label{eq:3Lexpansion}
\end{equation}
\end{lemma516}

\begin{proof}
By \eqref{eq:17}, there is an $R<\infty$ such that
\begin{equation}
 \varepsilon^2\frac{\|\nabla U(x)\|}{\|x\|}
 \geq\frac{16}{3}
 \qquad\text{whenever }\|x\|\geq R.
 \label{eq:gradientthreshold}
\end{equation}
The claim holds once $\|x_0\|$ is large enough. Choose $x_0$ such that $\|x_0\| \geq \max(R,\eta)$, meaning Lemma 5.15 applies and $\|x_0\|\geq R$.
The case $i=1$ then follows from Lemma 5.15.  Now suppose, for some $i\in\{1,\ldots,L-1\}$, that
$\|x_{i\varepsilon}\|\geq3\|x_{(i-1)\varepsilon}\|$. Taking norms in
\eqref{eq:recurrence} gives
\begin{align*}
 \frac{\|x_{(i+1)\varepsilon}\|}{\|x_{i\varepsilon}\|}
 &\geq
 \varepsilon^2
 \frac{\|\nabla U(x_{i\varepsilon})\|}{\|x_{i\varepsilon}\|}
 -2-
 \frac{\|x_{(i-1)\varepsilon}\|}{\|x_{i\varepsilon}\|}\\
 &\geq
 \varepsilon^2
 \frac{\|\nabla U(x_{i\varepsilon})\|}{\|x_{i\varepsilon}\|}
 -2-\frac13.
\end{align*}
The preceding steps of the induction imply
$\|x_{i\varepsilon}\|\geq\|x_0\|\geq R$, and therefore
\eqref{eq:gradientthreshold} gives
\[
 \frac{\|x_{(i+1)\varepsilon}\|}{\|x_{i\varepsilon}\|}
 \geq\frac{16}{3}-\frac{7}{3}=3.
\]
This proves \eqref{eq:stepwiseexpansion} by induction, and
\eqref{eq:3Lexpansion} follows immediately.
\end{proof}

We also retain condition (19): for the fixed number $L$ of leapfrog steps,
\begin{equation}
 \lim_{\substack{\|x\|\to\infty\\
                  \|y\|\geq2^L\|x\|}}
 \left(U(x)-U(y)-\frac12\|x\|^2\right)=-\infty.
 \tag{19}
 \label{eq:19}
\end{equation}

\begin{lemma517}
If \eqref{eq:17}, \eqref{eq:18} and \eqref{eq:19} hold, then for any
$\delta<1$ it holds that
\[
 \lim_{\substack{\|x_0\|\to\infty\\
                  \|p_0\|\leq\|x_0\|^\delta}}
 \alpha(x_0,x_{L\varepsilon})=0.
\]
\end{lemma517}

\begin{proof}
Take $\|x_0\|$ sufficiently large that Lemma 5.16 applies and that
$3\|x_{i\varepsilon}\|\geq C$ for $i=0,\ldots,L-1$. Using
\eqref{eq:stepwiseexpansion} and the corrected condition (18) then gives for $i \in \{0,...,L-1\}$
\begin{equation}
 \|\nabla U(x_{i\varepsilon})\|
 \leq3^{-(L-i)}\|\nabla U(x_{L\varepsilon})\|.
 \label{eq:gradientgeometric}
\end{equation}
Now, recall from equation (11) of \cite{Livingstone2019} that
\[
 p_{L\varepsilon}
 =p_0-\frac{\varepsilon}{2}\nabla U(x_0)
 -\varepsilon\sum_{i=1}^{L-1}\nabla U(x_{i\varepsilon})
 -\frac{\varepsilon}{2}\nabla U(x_{L\varepsilon}).
\]
Using \eqref{eq:gradientgeometric} and the reverse triangle inequality,
\begin{align*}
 \|p_{L\varepsilon}\|
 &\geq
 \frac{\varepsilon}{2}\|\nabla U(x_{L\varepsilon})\|
 -\varepsilon\sum_{i=1}^{L-1}\|\nabla U(x_{i\varepsilon})\|
 -\frac{\varepsilon}{2}\|\nabla U(x_0)\|-\|p_0\|\\
 &\geq
 \frac{\varepsilon}{2}\|\nabla U(x_{L\varepsilon})\|
 -\varepsilon\sum_{i=1}^{L}\left(\frac13\right)^i
 \|\nabla U(x_{L\varepsilon})\|-\|p_0\|\\
 &=
 \frac{\varepsilon}{2}
 \left\{1-2\sum_{i=1}^{L}\left(\frac13\right)^i\right\}
 \|\nabla U(x_{L\varepsilon})\|-\|p_0\|\\
 &=
 \frac{\varepsilon}{2\cdot3^L}
 \|\nabla U(x_{L\varepsilon})\|-\|p_0\|.
\end{align*}
Set $\gamma_L=3^{-L}$. Since \eqref{eq:17} holds and
$\|x_{L\varepsilon}\|\geq3^L\|x_0\|$, for all sufficiently large
$\|x_0\|$ the right-hand side above is positive. Squaring then gives
\begin{align*}
 \|p_{L\varepsilon}\|^2
 &\geq
 \left(
 \frac{\varepsilon\gamma_L}{2}\|\nabla U(x_{L\varepsilon})\|
 -\|p_0\|
 \right)^2\\
 &=
 \frac{\varepsilon^2\gamma_L^2}{4}
 \|\nabla U(x_{L\varepsilon})\|^2
 +\|p_0\|^2
 -\varepsilon\gamma_L
 \|\nabla U(x_{L\varepsilon})\|\|p_0\|.
\end{align*}
It follows that
\begin{equation}
 \|p_0\|^2-\|p_{L\varepsilon}\|^2
 \leq
 \varepsilon\gamma_L\|\nabla U(x_{L\varepsilon})\|
 \left(
 \|p_0\|-\frac{\varepsilon\gamma_L}{4}
 \|\nabla U(x_{L\varepsilon})\|
 \right).
 \label{eq:kineticbound}
\end{equation}
Again using \eqref{eq:17} and \eqref{eq:3Lexpansion}, uniformly over
$\|p_0\|\leq\|x_0\|^\delta$ we may take $\|x_0\|$ sufficiently large that
\[
 \frac{\varepsilon\gamma_L}{4}
 \|\nabla U(x_{L\varepsilon})\|\geq2\|x_0\|,
 \qquad
 \|p_0\|\leq\|x_0\|.
\]
The right-hand side of \eqref{eq:kineticbound} is then at most
$-\|x_0\|^2$, which implies
\[
 U(x_0)-U(x_{L\varepsilon})
 +\frac12\|p_0\|^2-\frac12\|p_{L\varepsilon}\|^2
\leq
 U(x_0)-U(x_{L\varepsilon})-\frac12\|x_0\|^2.
\]
By \eqref{eq:3Lexpansion},
$\|x_{L\varepsilon}\|\geq3^L\|x_0\|\geq2^L\|x_0\|$, so the right-hand
side tends to $-\infty$ by \eqref{eq:19}. Recalling equation (10) of
\cite{Livingstone2019}, this proves
$\alpha(x_0,x_{L\varepsilon})\to0$ uniformly over
$\|p_0\|\leq\|x_0\|^\delta$.
\end{proof}

\begin{theorem514}
If \eqref{eq:17}, \eqref{eq:18} and \eqref{eq:19} hold, then the Hamiltonian
Monte Carlo method with fixed integration time $T=L\varepsilon$ does not
produce a geometrically ergodic Markov chain for any choice $T>0$.
\end{theorem514}

\begin{proof}
Fix any $\delta<1$. From Lemma 5.17,
\[
 \sup_{\|p_0\|\leq\|x_0\|^\delta}
 \alpha(x_0,x_{L\varepsilon})\longrightarrow0
 \qquad\text{as }\|x_0\|\to\infty.
\]
Moreover,
\[
 \mu_G\!\left(\left\{p_0:\|p_0\|>\|x_0\|^\delta\right\}\right)
 \longrightarrow0.
\]
Therefore, writing $r(x_0)$ for the rejection probability as in Definition
3.2 of \cite{Livingstone2019},
\begin{align*}
 1-r(x_0)
 &=\int\alpha(x_0,x_{L\varepsilon})\,\mu_G(dp_0)\\
 &\leq
 \sup_{\|p_0\|\leq\|x_0\|^\delta}
 \alpha(x_0,x_{L\varepsilon})
 +\mu_G\!\left(\left\{p_0:\|p_0\|>\|x_0\|^\delta\right\}\right)
 \longrightarrow0.
\end{align*}
This implies that the rejection probability $r(x_0)\to1$ as $\|x_0\|\to\infty$, and Proposition 3.5 of
\cite{Livingstone2019} then implies that the chain is not geometrically ergodic.
\end{proof}

When the number of leapfrog steps is drawn from a position-independent
distribution $\mathcal{L}(\cdot)$ as in Assumption A1 of \cite{Livingstone2019}, the
same conclusion follows provided the above argument is applied conditional on
each fixed $L$. For each fixed $L$ the conditional acceptance probability
tends to zero, while it is bounded by one. Since the law of $L$ does not
depend on $x_0$, dominated convergence with respect to $\mathcal{L}(\cdot)$ shows that
the unconditional acceptance probability also tends to zero.

\section{Correction to Corollary 2.3}

On p.~3127, the statement that Corollary 2.3(ii) is a direct consequence of
Theorem 2.2 requires the following verification in the case $\beta>2$.

\begin{corollary23}
For the exponential family $E(\beta,\alpha)$ of \cite{Livingstone2019}, the
Hamiltonian Monte Carlo method does not produce a geometrically ergodic
Markov chain for every $\beta>2$ under Assumption A1.
\end{corollary23}

\begin{proof}
For some $M<\infty$, when $|x|> M$ we have by assumption that
$\pi(x)\propto\exp(-\alpha|x|^\beta)$,
so
\[
 \nabla U(x)=\alpha\beta\operatorname{sgn}(x)|x|^{\beta-1}.
\]
When $\beta>2$,
\[
 \frac{\|\nabla U(x)\|}{\|x\|}
 =\alpha\beta|x|^{\beta-2}\longrightarrow\infty,
\]
which is condition \eqref{eq:17}.

Choose $C>3M$. If $|y|\geq3|x|\geq C$, then both $x$ and $y$ are in the tail
region and
\[
 \frac{\|\nabla U(y)\|}{\|\nabla U(x)\|}
 =\left(\frac{|y|}{|x|}\right)^{\beta-1}
 \geq3^{\beta-1}>3.
\]
The corrected condition (18) therefore holds for every $\beta>2$.

Finally, for each fixed $L\geq1$, if $|y|\geq2^L|x|$ and $|x|$ is
sufficiently large, then
\begin{align*}
 U(x)-U(y)-\frac12|x|^2
 &\leq
 -\alpha\bigl(2^{L\beta}-1\bigr)|x|^\beta
 -\frac12|x|^2
 \longrightarrow-\infty,
\end{align*}
meaning condition \eqref{eq:19} holds for every fixed $L$. The corrected
Theorem 5.14 applies conditional on each fixed $L$, and dominated convergence
with respect to the position-independent distribution $\mathcal{L}(\cdot)$ gives the
corresponding result under Assumption A1. The Hamiltonian Monte Carlo
method therefore does not produce a geometrically ergodic Markov chain for every
$\beta>2$.
\end{proof}

The $\beta<1$ part of Corollary 2.3(ii), and all parts of Corollary 2.3(i),
are unaffected by this correction. Therefore Corollary 2.3 remains valid as
stated. No other changes to the results of \cite{Livingstone2019} are
required.

%% Optional acknowledgement, if appropriate:
%% \begin{acks}[Acknowledgments]
%% The authors thank [NAME] for bringing the issue corrected here to their attention.
%% \end{acks}

\begin{thebibliography}{}

\bibitem[\protect\citeauthoryear{Livingstone et al.}{2019}]{Livingstone2019}
Livingstone, S., Betancourt, M., Byrne, S. and Girolami, M. (2019).
On the geometric ergodicity of Hamiltonian Monte Carlo.
\textit{Bernoulli} \textbf{25} 3109--3138.
\href{https://doi.org/10.3150/18-BEJ1083}{doi:10.3150/18-BEJ1083}.

\bibitem{gruffaz2026theoretical}
Gruffaz, S., Kim, K., Guehtar, F., Duval-Decaix, H. and Trautmann, P. (2026).
A theoretical comparison of No-U-Turn Sampler variants: Necessary and sufficient
convergence conditions and mixing time analysis under Gaussian targets.
Preprint. Available at arXiv:2603.18640.

\end{thebibliography}

%\end{document}
